\section{微分方程}
\idx{阶}{jie}

微分方程中所含导数的最高阶数,称为该方程的\textbf{阶}(order).
例如 $f'+af=b$ 是 $1$ 阶,
$f''+p(x)f'+q(x)f=\lambda$ 是 $2$ 阶,
$\partial_x f+\partial_{yy}f+\partial_{xy}f=0$ 是 $2$ 阶,
$f^{(n)}-x f f^{(n-1)}=0$ 是 $n$ 阶.
有且仅有一个自变量的称为\textbf{常微分方程};
自变量多于一个的称为\textbf{偏微分方程}.
前两个例子是常微分方程,第三个是偏微分方程.
本讲要讨论的波动、热传导和稳定场方程都是偏微分方程.
书中典型的例子:贝塞尔方程是二阶线性齐次常微分方程,系数里的 $x^2$ 只依赖自变量,并不破坏线性;
拉普拉斯方程是二阶线性齐次偏微分方程,泊松方程右端换成已知函数 $f$ 后是线性非齐次;
热传导方程对时间是一阶,无源时线性齐次,加上已知热源后线性非齐次.
对时间的阶数决定初始条件要给几项.

方程对未知函数及其各阶导数总体为一次、系数只依赖自变量时,称为\textbf{线性}方程;否则为\textbf{非线性}方程.
非线性方程若只对最高阶导数为一次、系数可以含未知函数及其较低阶导数,称为\textbf{拟线性}方程,例如 $u_t+u u_x=0$.
最高阶导数本身不是一次的,称为\textbf{完全非线性}方程,例如 $(u_x)^2+(u_y)^2=u$.
本讲要导出的三条方程都是线性的.
右端不含未知函数及其导数时,称为齐次方程,否则为非齐次方程.
线性不一定齐次:右端可以是已知的外力或源.
边界条件和初始条件后面也这样区分.

物理规律讨论的是物理量在空间中的分布和在时间上的变化.
它反映的是某个物理量在邻近地点和邻近时刻之间的联系.
导出方程时,先确定所研究的物理量,
从系统中划出一个小部分,
按物理规律分析邻近部分与这个小部分的相互作用(抓住主要作用,略去次要因素),
以及这种作用在一个短时间段里怎样影响所研究的量,
再把这种影响写成算式并加以简化.
下面按这个步骤导出三条典型方程,然后再讨论泛定方程和定解问题.
